\section{Conclusion}
\label{sec:conclusion}

\subsection{Summary}

This paper introduced \ps{}$(p, T)$, a generalisation of \cs{} that makes
the compactness percentile an explicit statutory parameter.
By running $T$ bisection plans and returning the plan at rank $\floor{p
\times T}$, \ps{} allows the \dia{} to specify exactly which position in
the compactness distribution should be selected.

The central empirical finding is that partisan outcomes are largely
insensitive to the choice of $p$.
For six key states (NC, WI, GA, PA, TX, CA), projected Democratic seat counts
are identical across $p \in \{0.0, 0.25, 0.5, 0.75, 1.0\}$ with $T = 101$
plans.
The 0.5-seat bound is not violated for any state.
This confirms that compactness is determined by geographic structure, not by
search strategy---consistent with B.7 and B.17.

\subsection{The Legal Contribution}

The legal contribution of \ps{} is the separation of the \emph{algorithmic
posture} from the \emph{partisan outcome}.
Before \ps{}, the compactness extremum was the only option: \cs{} always
returned the minimum-edge-cut plan.
After \ps{}, the compactness percentile is a visible, debatable, statutory
parameter---like the population tolerance or the county edge weight.

This transparency is valuable even if the choice does not matter empirically.
A statute that specifies $p = 0.0$ and provides the legal rationale (no
neutral plan is more compact) is more robust to challenge than a statute
that produces the same plan without explaining why.
The insensitivity finding is itself part of the statutory record: the choice
of $p$ is not a partisan choice, so any value of $p$ is equally defensible
on neutrality grounds.

\subsection{Structure Is the Legally Relevant Choice}

The \ps{} framework confirms the portfolio's central thesis: algorithm
\emph{structure} (\ar{} vs. unweighted bisection, GeoSection vs. plain
bisection) is the legally relevant choice.
The search percentile is not.
A legislature that specifies a bisection algorithm has made the legally
consequential decision.
The percentile is a fine-tuning choice with predictable and auditable
compactness effects and no partisan consequence.

This finding strengthens the DIA's constitutional argument.
The statute specifies an algorithm structure (\ar{} + GeoSection) and a
compactness posture ($p = 0.0$).
Both choices are justified on non-partisan grounds.
The partition-space insensitivity of B.17 and the percentile insensitivity
of H.0 together establish that no neutral party could produce a different
partisan outcome by running any variant of the specified algorithm.

\subsection{Future Work}

Three open questions remain.

First, does the insensitivity result hold for $T < 101$?
With fewer plans, the distribution is less well-sampled and rank $r =
\floor{p \times T}$ is a cruder quantile estimate.
A sensitivity analysis over $T \in \{11, 21, 51, 101, 201\}$ would
bound the minimum $T$ required for stable percentile selection.

Second, does \ts{}(50th, \recom{}) produce the same partisan outcomes as
\ps{}$(0.5, T)$?
If so, the two postures are legally equivalent and the simpler \ps{}
implementation should be preferred.
If not, the \recom{} median lands in a genuinely different part of the plan
space, and the legal difference between bisection-family representativeness
and full-ensemble representativeness becomes substantive.

Third, the \ps{} framework extends naturally to multi-constraint variants
(VrASection, AreaSection).
Whether percentile insensitivity holds when additional constraints are added
to the bisection objective is an open empirical question.
